\songtitle{The Jingle Before the Words}
\setin{1987-12-17}

\tag*{2026}\tag{folk-rock}\tag{autobiographical}\label{sec:the-jingle-before-the-words}
\begin{notes}
Lyrics by Suno based on a Claude-produced excerpt of \emph{A Life in World Events}, the memoir of Carlos Thompson -- the Caracol news jingle that made a living room go still, from the fall of a regime in Haiti to 17 December 1987, driving home from a friends' novena when the radio announced the murder of Guillermo Cano. Opens and closes with lines in Spanish. Retrieved from Carlos Th's Suno page (V6-MINI, 13 November 2026).
\end{notes}
\begin{style}
1980s folk-rock ballad; male baritone with clear melodic delivery, an urgent spoken intro turning into crooning legato phrasing and doubled hooks; close-miked vocal with plate reverb; mid-tempo, melancolic pulse in D-minor; strummed acoustic guitar, acoustic drum kit and sustained electric guitar; fuller final chorus.
\end{style}
\begin{lyrics}
\songpart{Intro - spoken}
«¡Cuando la noticia se produce, Caracol se la comunica!»
«!Extra!\\
Ofrecemos en Caracol un boletín extraordinario de última hora»

\songpart{Verse 1 - sung, male singer}
By fifteen, I knew the sound\\
That made a living room go still\\
A jingle crossed the radio\\
And every voice obeyed its will\\
In Colombia, the words would come:\\
“Cuando la noticia se produce,\\
Caracol se la comunica.”\\
A hand raised gently in the air\\
Before the facts were anywhere

\songpart{Pre-Chorus}
I’d heard it when a far-off rule\\
Was falling down in Haiti\\
The tune said, something has occurred\\
And silence found the family

\songpart{Chorus}
The jingle came before the words\\
Before the night could tell us why\\
It turned the road to something else\\
Beneath a December sky\\
It never said that all was well\\
It never promised what would be\\
Only that the world had changed\\
And news was on its way to me

\songpart{Verse 2 - sung, male singer}
December seventeenth, ’87\\
We’d left our friends’ novena\\
Nine nights of prayer before the feast\\
A visit with our friends, I’d say\\
We drove toward home; the radio spoke\\
That same familiar little phrase\\
Then Guillermo Cano had died\\
Outside the paper where he stayed

\songpart{Pre-Chorus}
He’d challenged Medellín’s cartel\\
In print, in public, year by year\\
At fifteen I had read so little\\
Still I understood the fear

\songpart{Chorus}
The jingle came before the words\\
Before the night could tell us why\\
It turned the road to something else\\
Beneath a December sky\\
It never said that all was well\\
It never promised what would be\\
Only that the world had changed\\
And news was on its way to me

\songpart{Bridge - sung, male singer}
I can’t recall the rest of that drive\\
The turns, the houses passing by\\
I kept the sound, the sudden weight\\
That settled in before the name\\
And after Cano, violence grew\\
The headlines darkened, one by one\\
A song of news, a signal first\\
Then all the words we couldn’t outrun

\annotation{Spoken News Call}\\
Buenas noches. Guillermo Cano Isaza, director of El Espectador, has been assassinated outside the newspaper’s offices. He had publicly opposed the Medellín cartel.

\songpart{Final Chorus}
The jingle came before the words\\
And somehow I already knew\\
A room can change before a sentence\\
A road can split in two\\
It never said that all was well\\
It never promised what would be\\
Only that the world had changed\\
And news was on its way to me

\songpart{Outro}
“Cuando la noticia se produce,\\
Caracol se la comunica.”\\
The jingle never promised good news\\
Only that something had happened
\end{lyrics}

